% =====================================================================
% Section 3.2: Design and System Plans
% Reformatted per panel recommendation: numbered-list algorithm style,
% consolidated from 12 to 7 algorithms.
% =====================================================================

% Section heading is provided by the parent file (Chapter 3, SOP #2).
% This file begins directly with the introductory paragraph.

Section~3.2 records the design routines behind the implemented
\sysname{} prototype. Section~\ref{sec:architecture} (\S2.8.3)
identifies the major layers and responsibilities; this section moves
one level closer to the code by stating each routine's inputs, outputs,
equations, parameter values, and numbered steps. The discussion follows
the detection-to-ticket flow in Figure~\ref{fig:detection-ticket-pipeline}
(Figure~2.2): camera-frame filtering, queue-zone confirmation,
tracker maintenance, re-identification, queue-state updates, prediction,
and ticket output.

\subsection{Notation}
\label{subsec:notation}

The specifications below adopt the following notation. Queueing-theoretic
symbols ($\lambda$, $\mu$, $c$, $\rho$, $W$) retain the meanings introduced
in Section~\ref{sec:theoretical} (\S1.5). Queue length is denoted $Q$
throughout this section; $Q$ is identical to $L$ in
Equation~\ref{eq:littles_law}. Two distinct smoothing factors appear:
$\alpha_{\mathrm{H}}$ and $\beta_{\mathrm{H}}$ for Holt's
double-exponential smoothing of queue counts (\S\ref{subsec:prediction}),
and $\alpha_{\mathrm{s}}$ for the per-frame exponential moving average
applied to bounding-box coordinates (\S\ref{subsec:cv-layer}).
Thresholds are denoted $\tau$ with a descriptive subscript
($\tau_{\mathrm{match}}$ for the absolute face-match threshold,
$\tau_{\mathrm{margin}}$ for the required lead over the runner-up,
$\tau_{\mathrm{det}}$ for the face-detector confidence gate,
$\tau_{\mathrm{iou}}$ for the duplicate-suppression IoU gate).
The intersection-over-union of two axis-aligned bounding boxes $A$ and
$B$ is

\begin{equation}
\mathrm{IoU}(A, B) \;=\; \frac{|A \cap B|}{|A \cup B|},
\label{eq:iou}
\end{equation}

\noindent used by the duplicate-suppression step (\S\ref{subsec:dedup}),
the track-identifier remapping step (\S\ref{subsec:remap}), and the
non-maximum suppression step internal to the tracker
(\S\ref{subsec:cv-layer}).

% =====================================================================
\subsection{Computer Vision Layer}
\label{subsec:cv-layer}

The Computer Vision Layer runs as a separate Python process
(\texttt{app/detector.py}) on the camera-connected machine. Each
captured frame is fed to a YOLOv8n detector~\cite{jocher2023}
configured for single-class person detection. Multi-object tracking is
performed by ByteTrack~\cite{zhang2022bytetrack}, the default tracker
when the Ultralytics interface is invoked with \texttt{persist=True};
ByteTrack assigns a persistent track identifier $\mathrm{tid}$ to each
detected person across frames. Raw YOLO output is then passed through
five quality gates before any detection is forwarded to the Business
Logic Layer, and each accepted bounding box is subsequently smoothed
with an exponential moving average to remove per-frame jitter.

\subsubsection{Detection quality gates}

For a bounding box with corners $(x_1, y_1, x_2, y_2)$ in a frame of
size $W_f \times H_f$, Gate~1 requires sufficient area:

\begin{equation}
A \;=\; (x_2 - x_1)(y_2 - y_1) \;\geq\; A_{\min}.
\label{eq:gate1}
\end{equation}

\noindent Gate~2 caps the fraction of the frame a single box may occupy:

\begin{equation}
f \;=\; \frac{A}{W_f \, H_f} \;\leq\; f_{\max}.
\end{equation}

\noindent Gate~3 enforces a portrait-orientation aspect ratio:

\begin{equation}
r \;=\; \frac{y_2 - y_1}{x_2 - x_1} \;\geq\; r_{\min}.
\end{equation}

\noindent Gate~4 requires evidence of motion:

\begin{equation}
M \;=\; \big|\{p : |g_t(p) - g_{t-1}(p)| > 25\}\big| \;\geq\; M_{\min}
\label{eq:motion}
\end{equation}

\noindent or, as an override for static high-confidence detections,
$\text{conf} \geq c_{\mathrm{byp}}$. Gate~5 is standard
non-maximum suppression applied internally by ByteTrack. Only
detections that pass all five gates reach the queue tracker.
Parameter values are listed in Table~\ref{tab:cv-params}.

\subsubsection{Bounding-box smoothing}

Each of the four bounding-box coordinates is smoothed with an
exponential moving average before being forwarded:

\begin{equation}
\hat{b}_t \;=\; \alpha_{\mathrm{s}}\, b_t \;+\; (1 - \alpha_{\mathrm{s}})\, \hat{b}_{t-1},
\label{eq:ema-bbox}
\end{equation}

\noindent applied independently to $x_1$, $y_1$, $x_2$, and $y_2$.
Algorithm~1 below formalises the complete detection processing routine,
combining all five quality gates with bounding-box smoothing.

\vspace{0.5em}
\noindent\rule{\linewidth}{0.8pt}

\noindent\textbf{\underline{Algorithm 1: Detection Quality Filtering}}

\noindent\textbf{Input:} \textit{currentFrame} $F_t$,
\textit{previousGreyFrame} $G_{t-1}$, \textit{rawDetections} $D$ from YOLOv8n\\
\noindent\textbf{Output:} \textit{filteredDetections} (quality-gated
detections with smoothed bounding boxes, ready for queue-zone
processing)

\begin{enumerate}
  \item \textbf{CONVERT} $F_t$ to greyscale to obtain $G_t$;
        \textbf{INITIALIZE} $\mathcal{D}_{\mathrm{ok}} \gets \emptyset$
  \item \textbf{FOR EACH} detection $d \in D$ with box
        $(x_1,y_1,x_2,y_2)$ and confidence $\text{conf}$:
  \item \quad \textbf{COMPUTE} $A = (x_2-x_1)(y_2-y_1)$,
        $f = A/(W_fH_f)$, $r = (y_2-y_1)/(x_2-x_1)$
  \item \quad \textbf{IF} $A < A_{\min}$ \textbf{OR} $f > f_{\max}$
        \textbf{OR} $r < r_{\min}$ \textbf{THEN} skip $d$
        (fails Gates 1--3)
  \item \quad \textbf{COMPUTE} motion energy $M$ as the pixel count
        where $|G_t(p) - G_{t-1}(p)| > 25$
  \item \quad \textbf{IF} $M < M_{\min}$ \textbf{AND}
        $\text{conf} < c_{\mathrm{byp}}$ \textbf{THEN} skip $d$
        (fails Gate 4)
  \item \quad \textbf{APPEND} $d$ to $\mathcal{D}_{\mathrm{ok}}$
  \item \textbf{APPLY} ByteTrack NMS to $\mathcal{D}_{\mathrm{ok}}$
        (Gate 5)
  \item \textbf{FOR EACH} accepted detection \textbf{APPLY} EMA
        smoothing: $\hat{b}_t \gets \alpha_{\mathrm{s}}\,b_t +
        (1-\alpha_{\mathrm{s}})\,\hat{b}_{t-1}$ on each coordinate
        independently
  \item \textbf{RETURN} $\mathcal{D}_{\mathrm{ok}}$ with smoothed
        bounding boxes
\end{enumerate}
\vspace{0.3em}

\begin{table}[H]
\centering
\caption{Computer Vision Layer parameters.}
\label{tab:cv-params}
\begin{tabular}{lcl}
\hline
Symbol & Value & Description \\
\hline
$A_{\min}$       & $800\ \mathrm{px}^2$ & Minimum bounding-box area (Gate~1) \\
$f_{\max}$       & $0.85$               & Maximum frame coverage (Gate~2) \\
$r_{\min}$       & $0.60$               & Minimum portrait aspect ratio (Gate~3) \\
$M_{\min}$       & $8\ \mathrm{px}$     & Minimum motion energy (Gate~4) \\
$c_{\mathrm{byp}}$ & $0.70$             & Static-detection confidence bypass (Gate~4) \\
$\alpha_{\mathrm{s}}$ & $0.45$          & Bounding-box EMA smoothing factor \\
\hline
\end{tabular}
\end{table}

% =====================================================================
\subsection{Queue Zone Membership and Candidate Confirmation}
\label{subsec:candidate}

A configurable rectangular region of interest, the queue zone, is
defined by four pixel coordinates $(z_{x_1}, z_{y_1}, z_{x_2}, z_{y_2})$
on the camera frame. A detected person is considered \emph{inside} the
zone when the centroid of their bounding box,
$(c_x, c_y) = \big((x_1+x_2)/2,\ (y_1+y_2)/2\big)$, falls within
the rectangle:

\begin{equation}
\text{inside} \;\Longleftrightarrow\; z_{x_1} \leq c_x \leq z_{x_2}
\;\wedge\; z_{y_1} \leq c_y \leq z_{y_2}.
\label{eq:zone}
\end{equation}

\noindent A person entering the queue zone is not assigned a queue
number immediately. Each new track identifier is buffered in a
candidate dictionary $C$ and must satisfy persistence and motion
conditions before a number is issued. Persistence requires the same
$\mathrm{tid}$ remain inside the zone for at least $N_{\mathrm{conf}}$
consecutive frames. Motion is assessed from the standard deviations
of the buffered centroid sequence:

\begin{equation}
\sigma_x = \mathrm{std}(c_{x,1},\ldots,c_{x,N_{\mathrm{conf}}}),
\qquad
\sigma_y = \mathrm{std}(c_{y,1},\ldots,c_{y,N_{\mathrm{conf}}}),
\label{eq:sigma}
\end{equation}

\begin{equation}
\text{motion ok} \;\Longleftrightarrow\;
\sigma_x + \sigma_y > \theta_{\mathrm{m}}
\;\vee\;
\overline{\text{conf}} \geq c_{\mathrm{byp}}.
\label{eq:motion-gate}
\end{equation}

\vspace{0.5em}
\noindent\rule{\linewidth}{0.8pt}

\noindent\textbf{\underline{Algorithm 2: Candidate Confirmation}}

\noindent\textbf{Input:} \textit{detection} $d$ with track ID
$\mathrm{tid}$ and confidence $\text{conf}$, \textit{candidateBuffer}
$C$, \textit{confirmationThreshold} $N_{\mathrm{conf}}$\\
\noindent\textbf{Output:} \textit{assignmentDecision} (issue queue
number, continue accumulating, or reject as static)

\begin{enumerate}
  \item \textbf{INCREMENT} $C[\mathrm{tid}].\text{count}$ by 1;
        \textbf{APPEND} $(c_x, c_y)$ and $\text{conf}$ to
        $C[\mathrm{tid}]$
  \item \textbf{IF} $C[\mathrm{tid}].\text{count} < N_{\mathrm{conf}}$
        \textbf{THEN RETURN} (continue accumulating)
  \item \textbf{COMPUTE} $\sigma_x$, $\sigma_y$ over buffered
        centroid sequence in $C[\mathrm{tid}]$
  \item \textbf{COMPUTE} $\overline{\text{conf}} \gets
        \mathrm{mean}(C[\mathrm{tid}].\text{confs})$
  \item \textbf{IF} $\sigma_x + \sigma_y > \theta_{\mathrm{m}}$
        \textbf{OR} $\overline{\text{conf}} \geq c_{\mathrm{byp}}$
        \textbf{THEN:}
  \item \quad \textbf{ASSIGN} next available queue number to
        $\mathrm{tid}$; \textbf{GENERATE} ticket payload
  \item \textbf{ELSE:}
  \item \quad \textbf{DELETE} $C[\mathrm{tid}]$ (static object
        rejected)
  \item \textbf{RETURN} assignment decision
\end{enumerate}
\vspace{0.3em}

\noindent At $N_{\mathrm{conf}} = 20$ frames and $12$--$15$~fps, a
person must remain in the zone for approximately $1.3$--$1.7$~s before
receiving a queue number. The motion threshold
$\theta_{\mathrm{m}} = 1.5\ \mathrm{px}$ rejects static mis-detections;
the confidence bypass at $0.70$ admits a momentarily stationary person
whose detection is unambiguous.

% =====================================================================
\subsection{Track-Identifier Remapping}
\label{subsec:remap}

\subsection{Duplicate Suppression}
\label{subsec:dedup}

ByteTrack occasionally assigns a new $\mathrm{tid}$ to a person who
briefly left and re-entered the camera frame (handled by Part~A
below), or assigns two distinct $\mathrm{tid}$ values to the same
physical person within the same frame (a \emph{ghost track}, handled
by Part~B). Both cases are resolved in Algorithm~3.

Remapping uses IoU (Equation~\ref{eq:iou}) and centroid Euclidean
distance:

\begin{equation}
d(A, B) \;=\; \sqrt{(c_x^A - c_x^B)^2 + (c_y^A - c_y^B)^2},
\label{eq:cdist}
\end{equation}

\noindent combined into a remap score with IoU prioritised:

\begin{equation}
s \;=\;
\begin{cases}
\mathrm{IoU}(A, B) + 1.0 & \text{if } \mathrm{IoU}(A, B) \geq \tau_{\mathrm{iou}}^{\mathrm{rm}}, \\[2pt]
1 - d(A, B)/d_{\max}     & \text{else if } d(A, B) \leq d_{\max}, \\[2pt]
\text{(no match)}        & \text{otherwise.}
\end{cases}
\label{eq:remap}
\end{equation}

Ghost-track merging uses:

\begin{equation}
\mathrm{dup}(p_i, p_j) \;\Longleftrightarrow\;
\mathrm{IoU}(p_i, p_j) > \tau_{\mathrm{iou}}^{\mathrm{dd}}
\;\vee\;
d(p_i, p_j) < f_{\mathrm{dd}} \cdot \overline{D}(p_i, p_j),
\label{eq:dedup}
\end{equation}

\noindent where $\overline{D}(p_i, p_j)$ is the mean diagonal of the
two bounding boxes, and the younger of the two entries is dropped. The
test is purely geometric. A person who has just been restored after an
absence carries 60 frames of immunity from this pass, so that somebody
returning to the queue is not immediately removed as a duplicate of
themselves. Parameter values are in Table~\ref{tab:dedup-params}.

\begin{table}[H]
\centering
\caption{Track-management parameters.}
\label{tab:dedup-params}
\begin{tabular}{lcl}
\hline
Symbol & Value & Description \\
\hline
$\tau_{\mathrm{iou}}^{\mathrm{rm}}$ & $0.10$ & Remap IoU threshold \\
$d_{\max}$                          & $180\ \mathrm{px}$ & Remap distance cap \\
$F_{\mathrm{rm}}$                   & $45\ \text{frames}$ & Remap absence window \\
$\tau_{\mathrm{iou}}^{\mathrm{dd}}$ & $0.40$ & Ghost-track merge IoU \\
$f_{\mathrm{dd}}$                   & $0.15$ & Centre-distance fraction \\
$I_0$                               & $60\ \text{frames}$ & Re-entry immunity \\
\hline
\end{tabular}
\end{table}

\vspace{0.5em}
\noindent\rule{\linewidth}{0.8pt}

\noindent\textbf{\underline{Algorithm 3: Track Management ---
Remapping and Duplicate Suppression}}

\noindent\textbf{Input:} \textit{incomingDetection} $d$ with new
track ID $\mathrm{tid}_{\mathrm{new}}$, \textit{activeQueue}
$\mathcal{Q}$, \textit{immunityTable} $I$\\
\noindent\textbf{Output:} \textit{correctedTrackID}; updated
$\mathcal{Q}$ with ghost tracks removed

\noindent\textit{Part A --- Track-Identifier Remapping:}
\begin{enumerate}
  \item \textbf{SET} $\mathrm{bestID} \gets \mathrm{tid}_{\mathrm{new}}$;
        $\mathrm{bestScore} \gets -\infty$
  \item \textbf{FOR EACH} active person $p \in \mathcal{Q}$:
  \item \quad \textbf{COMPUTE} IoU $u$ between $d.\mathrm{bbox}$ and
        $p.\mathrm{bbox}$; \textbf{COMPUTE} centroid distance $\delta$
  \item \quad \textbf{IF} $u \geq \tau_{\mathrm{iou}}^{\mathrm{rm}}$
        \textbf{THEN} $s \gets u + 1.0$
  \item \quad \textbf{ELSE IF} $\delta \leq d_{\max}$ \textbf{THEN}
        $s \gets 1 - \delta/d_{\max}$
  \item \quad \textbf{ELSE CONTINUE}
  \item \quad \textbf{IF} $s > \mathrm{bestScore}$ \textbf{THEN}
        \textbf{SET} $\mathrm{bestScore} \gets s$;
        $\mathrm{bestID} \gets p.\mathrm{tid}$
  \item \textbf{RETURN} $\mathrm{bestID}$
\end{enumerate}

\noindent\textit{Part B --- Per-Frame Duplicate Suppression:}
\begin{enumerate}
  \setcounter{enumi}{8}
  \item \textbf{FOR EACH} ordered pair $(p_i, p_j) \in \mathcal{Q}
        \times \mathcal{Q}$ with $i < j$:
  \item \quad \textbf{IF} $I[p_i] > 0$ \textbf{OR} $I[p_j] > 0$
        \textbf{THEN CONTINUE} (immunity active)
  \item \quad \textbf{IF} appearance signatures exist \textbf{AND}
        $\mathrm{sim}(p_i, p_j) < \tau_{\mathrm{bl}}$ \textbf{THEN
        CONTINUE} (visually distinct)
  \item \quad \textbf{IF} $\mathrm{dup}(p_i, p_j)$ holds
        \textbf{THEN} drop the higher queue number; retain the lower
  \item \textbf{DECREMENT} every positive entry of $I$ by 1
\end{enumerate}
\vspace{0.3em}

% =====================================================================
\subsection{Face-Based Identity Validation}
\label{subsec:reid}

This algorithm answers the question the pre-oral panel required the system
to answer: given a person standing in the queue zone, which enrolled
student is this, if any? It serves two purposes in the running system.
First, it identifies a registered student so that a queue number can be
issued without manual input, which is the subject of
Statement of the Problem~2. Second, it restores a queue person whose
track was lost, typically because they stepped out of the camera's view
and returned, recovering their original number instead of issuing a
second one.

The earlier design answered both questions from clothing colour. That
approach was replaced because a clothing signature identifies an
appearance, not a person: two students in similar uniforms are not
separable by it, the same student changes appearance between visits, and
nothing in it can confirm that the person at the camera is the account
holder requesting service. A face signature carries identity directly, and
the tiered arrangement of recency windows, spatial fallbacks and lookalike
guards that existed to compensate for a weak signal is no longer needed.

\subsubsection{Face embedding}
\label{subsubsec:signature}

A face detector locates the face within the upper region of each person's
bounding box, and is accepted only when its detection confidence reaches
$\tau_{\mathrm{det}} = 0.60$. The aligned face crop is passed to an
ArcFace recognition model, which maps it to a unit-length embedding

\begin{equation}
e \;=\; f(\mathbf{x}) \in \mathbb{R}^{512},
\qquad \lVert e \rVert_2 = 1,
\label{eq:embedding}
\end{equation}

\noindent where $\mathbf{x}$ is the face crop and $f$ the recognition
network. The mapping is one-way: $\mathbf{x}$ cannot be recovered from
$e$. Both models are supplied by the InsightFace
\texttt{buffalo\_s} pack and execute on the CPU, with the detector
operating at a $320 \times 320$ input size.

\subsubsection{Enrolment signature}
\label{subsubsec:enrolment}

A student enrols by submitting $n \in [1, 5]$ photographs at varied angle
and lighting. Each photograph that yields a confident face contributes an
embedding, and the stored signature is their normalised mean:

\begin{equation}
e^{\mathrm{enrol}} \;=\;
\frac{\frac{1}{n}\sum_{i=1}^{n} e_i}
     {\bigl\lVert \frac{1}{n}\sum_{i=1}^{n} e_i \bigr\rVert_2}.
\label{eq:enrol}
\end{equation}

\noindent Averaging across angles makes the stored signature more robust
than any single photograph. The photographs themselves are discarded once
the embedding is computed and are never written to storage.

\subsubsection{Comparison metric}
\label{subsubsec:psim}

A live embedding is compared with an enrolled one by cosine similarity,
which for unit-length vectors reduces to their inner product:

\begin{equation}
\mathrm{sim}(e_1, e_2) \;=\;
\frac{e_1 \cdot e_2}{\lVert e_1 \rVert_2 \, \lVert e_2 \rVert_2}
\;=\; e_1 \cdot e_2 \;\in\; [-1, 1].
\label{eq:cosine}
\end{equation}

\subsubsection{Threshold-and-margin decision rule}
\label{subsubsec:hybrid}

Let $\mathcal{E} = \{(s_i,\, e_i^{\mathrm{enrol}})\}$ be the enrolled
students eligible for matching, and let the similarities of a live
embedding $e$ against $\mathcal{E}$ be ranked so that $s^{(1)}$ is the
highest and $s^{(2)}$ the next. Acceptance requires \emph{both} of:

\begin{equation}
\text{accept} \;\Longleftrightarrow\;
\underbrace{s^{(1)} \geq \tau_{\mathrm{match}}}_{\text{resembles a known student}}
\;\wedge\;
\underbrace{s^{(1)} - s^{(2)} \geq \tau_{\mathrm{margin}}}_{\text{resembles one student distinctly}},
\label{eq:face-decision}
\end{equation}

\noindent with $\tau_{\mathrm{match}} = 0.30$ and
$\tau_{\mathrm{margin}} = 0.15$. When only one student is enrolled,
$s^{(2)}$ is defined as $-1$, so the margin condition is satisfied
trivially and the absolute threshold alone governs.

The second condition is the element original to this system. The first
establishes only that the face resembles a known student; it does not
establish \emph{which} student, and an absolute threshold alone will
accept a face that resembles two enrolled students almost equally. The
margin condition requires the best match to be better than the runner-up
by a stated amount, so an ambiguous face is refused rather than resolved
by taking the higher of two close scores. A refusal is not a failure of
the system but its intended behaviour: the person remains pending and a
member of staff resolves the case, which is preferable to issuing a
confident but unverified identification. This also removes the need for a
dedicated lookalike guard, since ambiguity is rejected by construction
rather than by a special case.

The same rule governs the open-set requirement. People who have not
enrolled are never identified, because a face absent from $\mathcal{E}$
produces neither a sufficient absolute score nor a sufficient margin.

\subsubsection{Queue-number issuance}
\label{subsubsec:minting}

An accepted match is what allows a queue number to be issued without
manual input. Three conditions govern the issuance itself.

\textbf{Consent precedes issuance.} Recognition alone does not place a
student in the queue. The student must have armed a join intent in the
mobile client beforehand; a registered student who is recognized without
one is classified as a bystander and issued nothing, so that somebody
merely walking past the camera, accompanying a friend or working nearby,
is never added to a queue they did not ask to join. The intent lapses
after a configurable interval and is consumed once it produces a number.

\textbf{One active entry per student.} Before a number is issued, the
system checks whether that student already holds a pending or waiting
entry. If so, nothing is issued. A student obtains a new number only
after staff mark the previous one done, which prevents a
recognized student from accumulating duplicate numbers by re-entering the
camera's view.

\textbf{A separate numeric range.} System-issued numbers are drawn from a
reserved range beginning at $N_{\mathrm{face}} = 5000$:

\begin{equation}
\mathrm{number} \;=\; \min\bigl\{\, n \in \mathbb{N} \;:\;
n \geq N_{\mathrm{face}} \;\wedge\; n \notin \mathcal{U} \,\bigr\},
\label{eq:mint}
\end{equation}

\noindent where $\mathcal{U}$ is the set of numbers already in use. The
service office's existing kiosk issues ordinary low numbers, so the two
sequences cannot collide. This is what allows automatic issuance for
registered students while leaving the established queue flow untouched,
as the pre-oral panel required. Walk-ins never reach this path: an
unenrolled face produces no accepted match, and walk-ins continue to
take a ticket from the existing kiosk, which runs separately from
\sysname{}.

\vspace{0.5em}
\noindent\rule{\linewidth}{0.8pt}

\noindent\textbf{\underline{Algorithm 4: Face-Based Identity Validation}}

\noindent\textbf{Input:} \textit{faceCrop} $\mathbf{x}$ with detector
confidence $c$, \textit{eligibleEnrolled} $\mathcal{E}$ (students not
already served today), \textit{usedNumbers}
$\mathcal{U}$\\
\noindent\textbf{Output:} an issued \textit{queueNumber},
\textsc{Bystander}, \textsc{AlreadyQueued}, or \textsc{Unrecognized}

\begin{enumerate}
  \item \textbf{IF} $c < \tau_{\mathrm{det}}$ \textbf{THEN RETURN}
        \textsc{Unrecognized} (no usable face in this frame)
  \item \textbf{COMPUTE} $e \gets f(\mathbf{x})$, the $512$-dimensional
        unit-length embedding
  \item \textbf{IF} $\mathcal{E} = \emptyset$ \textbf{THEN RETURN}
        \textsc{Unrecognized}
  \item \textbf{COMPUTE} $\mathrm{sim}(e,\, e_i^{\mathrm{enrol}})$ for
        every $(s_i,\, e_i^{\mathrm{enrol}}) \in \mathcal{E}$
  \item \textbf{RANK} the similarities; \textbf{SET} $s^{(1)}$ and
        $s^{(2)}$ as the two highest, with $s^{(2)} \gets -1$ if
        $|\mathcal{E}| = 1$, and $p^{(1)}$ as the best-matching student
  \item \textbf{SET} $\mathrm{margin} \gets s^{(1)} - s^{(2)}$
  \item \textbf{RECORD} the match event with $s^{(1)}$, $\mathrm{margin}$,
        and the accept-or-reject outcome
  \item \textbf{IF} $s^{(1)} < \tau_{\mathrm{match}}$
        \textbf{OR} $\mathrm{margin} < \tau_{\mathrm{margin}}$
        \textbf{THEN RETURN} \textsc{Unrecognized}
        (below threshold, or ambiguous between candidates)
  \item \textbf{IF} $p^{(1)}$ has not armed a join intent
        \textbf{THEN} mark the person a bystander and
        \textbf{RETURN} \textsc{Bystander} (recognition is not consent)
  \item \textbf{IF} $p^{(1)}$ already holds a pending or waiting entry
        \textbf{THEN RETURN} \textsc{AlreadyQueued}
  \item \textbf{SET} $\mathrm{number} \gets
        \min\{\, n \geq N_{\mathrm{face}} : n \notin \mathcal{U} \,\}$
  \item \textbf{LINK} $\mathrm{number}$ to $p^{(1)}$;
        \textbf{ADD} $\mathrm{number}$ to $\mathcal{U}$;
        \textbf{CONSUME} the join intent
  \item \textbf{RETURN} $\mathrm{number}$
\end{enumerate}
\vspace{0.3em}

\noindent The set $\mathcal{E}$ is filtered by the caller to exclude
students already served for the day, which performs the
role the served-blacklist held in the earlier design without retaining any
appearance data. Every evaluation, accepted or refused, is written to a
recognition audit table together with its score and margin, which is the
source of the accuracy results reported in Chapter~3.

Re-entry restoration, the second purpose named at the start of this
section, applies the same decision rule to a different candidate set. When
a queue person's track is lost and a new detection appears, the live
embedding is compared against the stored embeddings of queue persons
currently recorded as missing rather than against the enrolled roster, and
Equation~\ref{eq:face-decision} governs acceptance unchanged. A match
restores the person's original queue number; a refusal leaves them to be
treated as a new arrival. Because one rule serves both purposes, the
system has a single notion of what counts as a confident identification.

% =====================================================================
\subsection{Counter Assignment}
\label{subsec:counters}

After every frame, the active queue $\mathcal{Q}$ is sorted by queue
number in ascending order. Counter assignment is \emph{sticky}: once a
counter is allocated to a queue entry, the allocation persists until
that entry is marked done or a no-show. The assignment rule is

\begin{equation}
\text{counter}(p) \;=\;
\begin{cases}
\text{rank}(p) & \text{if } \text{rank}(p) \leq c
                 \text{ and } p \text{ unassigned}, \\
\text{previous allocation} & \text{if } p \text{ already assigned}, \\
\text{(none)} & \text{otherwise.}
\end{cases}
\label{eq:counter}
\end{equation}

\vspace{0.5em}
\noindent\rule{\linewidth}{0.8pt}

\noindent\textbf{\underline{Algorithm 5: Sticky Counter Assignment}}

\noindent\textbf{Input:} \textit{activeQueue} $\mathcal{Q}$,
\textit{numberOfCounters} $c$\\
\noindent\textbf{Output:} updated $\mathcal{Q}$ with counter
assignments

\begin{enumerate}
  \item \textbf{SORT} $\mathcal{Q}$ by queue number ascending;
        \textbf{ASSIGN} position rank $1, 2, \ldots$ to each entry
  \item \textbf{SET} $\mathrm{occupied} \gets
        \{p.\text{counter} : p \in \mathcal{Q},\;
        p.\text{counter} \neq \text{none}\}$
  \item \textbf{SET} $\mathrm{free} \gets \{1,\ldots,c\}
        \setminus \mathrm{occupied}$, sorted ascending
  \item \textbf{FOR EACH} $p \in \mathcal{Q}$ with
        $p.\text{counter} = \text{none}$:
  \item \quad \textbf{IF} $\mathrm{free}$ is non-empty
        \textbf{THEN SET} $p.\text{counter} \gets$
        pop smallest element of $\mathrm{free}$
\end{enumerate}
\vspace{0.3em}

% =====================================================================
\subsection{Wait-Time Prediction}
\label{subsec:prediction}

The Prediction Service consumes a rolling buffer of up to sixty
crowd-count samples and updates the wait-time estimates every frame.
The service rate $\mu = 1/\overline{\text{svc}}$ is supplied by a
background measurement routine that re-estimates the average service
time every $T_{\mathrm{svc}} = 300\ \mathrm{s}$ from the MySQL
transaction log. For completed-transaction timestamps
$t_1, t_2, \ldots, t_N$, inter-departure gaps in minutes are

\begin{equation}
g_i \;=\; \frac{t_i - t_{i-1}}{60}, \qquad i = 2, \ldots, N.
\label{eq:gap}
\end{equation}

\noindent Gaps outside $[0.1,\, 15]\ \mathrm{min}$ are discarded.
The mean retained gap is multiplied by $c$ to correct for parallel
service, then blended with the previous estimate:

\begin{equation}
\widehat{\mathrm{svc}}_{\text{new}} \;=\;
0.7\,\widehat{\mathrm{svc}} + 0.3\,\widehat{\mathrm{svc}}_{\text{prev}}.
\label{eq:svc-blend}
\end{equation}

\subsubsection{M/M/c baseline (current wait)}
\label{subsubsec:mmc-baseline}

With $\rho = \lambda/(c\mu)$, the current wait estimate is:

\begin{equation}
W_{\text{current}} \;=\;
\begin{cases}
\dfrac{Q}{\max(\lambda,\, 0.1)} + \dfrac{1}{\mu}
  & \text{if } \rho < 1, \\[10pt]
\dfrac{Q}{\mu \, c}
  & \text{if } \rho \geq 1.
\end{cases}
\label{eq:wcurrent}
\end{equation}

\subsubsection{Arrival rate}
\label{subsubsec:arrival-rate}

\begin{equation}
\lambda \;=\; \frac{Q_t - Q_{t-n}}{\Delta t},
\qquad n = \min(30,\, |\text{history}|).
\label{eq:arrival}
\end{equation}

\subsubsection{Short-term forecast ($\approx 5\ \mathrm{min}$)}
\label{subsubsec:short-forecast}

An OLS slope $\beta$ is fit to the most recent $W_{\mathrm{tr}} = 20$
samples. The projected queue count is:

\begin{equation}
\hat{Q}_{\text{short}} \;=\;
\max\!\big(0,\, Q + \beta \cdot \mathrm{fps} \cdot T_{\mathrm{sh}}\big).
\label{eq:qshort}
\end{equation}

\subsubsection{Medium-term forecast ($\approx 15\ \mathrm{min}$)}
\label{subsubsec:holt-forecast}

Holt's double-exponential smoothing updates level $L_t$ and trend
$T_t$:

\begin{align}
L_t &\;=\; \alpha_{\mathrm{H}}\,y_t +
(1-\alpha_{\mathrm{H}})(L_{t-1}+T_{t-1}),
\label{eq:holt-level} \\
T_t &\;=\; \beta_{\mathrm{H}}(L_t - L_{t-1}) +
(1-\beta_{\mathrm{H}})T_{t-1}.
\label{eq:holt-trend}
\end{align}

\noindent The damped forecast $h$ steps ahead is:

\begin{equation}
\hat{Q}_{t+h} \;=\; L_t \;+\; T_t \sum_{i=1}^{h} \varphi^i.
\label{eq:holt-forecast}
\end{equation}

\subsubsection{Long-term forecast ($\approx 30\ \mathrm{min}$)}
\label{subsubsec:long-forecast}

The growth ratio and mean-reversion projection are:

\begin{equation}
g \;=\;
\frac{Q_{\mathrm{recent}}}{\max(Q_{\mathrm{early}},\, 1)},
\label{eq:growth}
\end{equation}

\begin{equation}
\hat{Q}_{\text{long}} \;=\;
\max\!\Big(0,\;
Q_{\mathrm{recent}} \cdot (0.7\,g + 0.3)\Big).
\label{eq:qlong}
\end{equation}

\subsubsection{Hybrid AI predictor}
\label{subsubsec:hybrid-ai-prediction}

The deployed AI model uses the analytical outputs
$W_{\mathrm{current}}$, $W_{\mathrm{short}}$, $W_{\mathrm{medium}}$,
and $W_{\mathrm{long}}$ as features for a Gradient Boosting regressor,
calibrated at runtime with current service speed, yesterday's same-hour
wait, and recent prediction error.

\vspace{0.5em}
\noindent\rule{\linewidth}{0.8pt}

\noindent\textbf{\underline{Algorithm 6: Multi-Horizon Wait-Time
Prediction}}

\noindent\textbf{Input:} \textit{countHistory} $H$,
\textit{currentQueueLength} $Q$, \textit{activeCounters} $c$,
\textit{serviceRate} $\mu$ (from dynamic measurement),
\textit{trainedModel} $G$\\
\noindent\textbf{Output:} \textit{currentWait} $W_{\mathrm{current}}$,
\textit{shortWait} $W_{\mathrm{short}}$,
\textit{mediumWait} $W_{\mathrm{medium}}$,
\textit{longWait} $W_{\mathrm{long}}$,
\textit{finalWait} $W_{\mathrm{final}}$

\begin{enumerate}
  \item \textbf{APPEND} current $Q$ and timestamp to $H$
  \item \textbf{COMPUTE} $\lambda$ from recent count differences
        using Equation~\ref{eq:arrival};
        \textbf{COMPUTE} $\rho \gets \lambda/(c\mu)$
  \item \textbf{COMPUTE} $W_{\mathrm{current}}$ using M/M/c formula
        (Equation~\ref{eq:wcurrent})
  \item \textbf{FIT} OLS slope $\beta$ to the latest
        $W_{\mathrm{tr}}$ samples;
        \textbf{COMPUTE} $\hat{Q}_{\text{short}} \gets
        \max(0, Q + \beta \cdot \mathrm{fps} \cdot T_{\mathrm{sh}})$
  \item \textbf{SET} $W_{\mathrm{short}} \gets$ M/M/c wait from
        $\hat{Q}_{\text{short}}$
  \item \textbf{UPDATE} Holt level $L_t$ and trend $T_t$ using
        $\alpha_{\mathrm{H}}$, $\beta_{\mathrm{H}}$
        (Equations~\ref{eq:holt-level}--\ref{eq:holt-trend})
  \item \textbf{COMPUTE} $\hat{Q}_{t+h} \gets L_t + T_t
        \sum_{i=1}^{h}\varphi^i$;
        \textbf{SET} $W_{\mathrm{medium}} \gets$ M/M/c wait from
        $\hat{Q}_{t+h}$
  \item \textbf{SPLIT} $H$ into thirds;
        \textbf{COMPUTE} growth ratio $g \gets Q_{\mathrm{recent}} /
        \max(Q_{\mathrm{early}}, 1)$
  \item \textbf{COMPUTE} $\hat{Q}_{\text{long}} \gets
        \max\!\big(0,\; Q_{\mathrm{recent}}(0.7g + 0.3)\big)$
  \item \textbf{SET} $W_{\mathrm{long}} \gets
        \min\!\big(W_{\max}^{\mathrm{long}},\;
        \text{M/M/c wait from } \hat{Q}_{\text{long}}\big)$
  \item \textbf{BUILD} feature vector $x$ from queue-state features,
        recent-demand features, academic-period indicators, and
        $\{W_{\mathrm{current}}, W_{\mathrm{short}},
        W_{\mathrm{medium}}, W_{\mathrm{long}}\}$
  \item \textbf{COMPUTE} $W_{\mathrm{AI}} \gets G(x)$;
        \textbf{APPLY} recent error correction $e_{\mathrm{recent}}$:
        $W_{\mathrm{AI,corr}} \gets \max(0, W_{\mathrm{AI}} +
        e_{\mathrm{recent}})$
  \item \textbf{COMPUTE} runtime baseline $W_{\mathrm{base}}$ from
        current service speed and yesterday's same-hour average
  \item \textbf{SET} $\gamma \gets
        \min(\text{completed\_today}/10,\, 1)$
  \item \textbf{COMPUTE} $W_{\mathrm{final}} \gets
        0.6\,W_{\mathrm{AI,corr}} + 0.4\,W_{\mathrm{base}}$
  \item \textbf{RETURN}
        $(W_{\mathrm{current}}, W_{\mathrm{short}},
        W_{\mathrm{medium}}, W_{\mathrm{long}},
        W_{\mathrm{final}}, \rho, \lambda)$
\end{enumerate}
\vspace{0.3em}

The deployed parameter values are $\alpha_{\mathrm{H}} = 0.4$,
$\beta_{\mathrm{H}} = 0.2$, $\varphi = 0.85$, $h = 20$, and
$W_{\max}^{\mathrm{long}} = 60\ \mathrm{min}$.

% =====================================================================
\subsection{Dynamic Service-Time Measurement}
\label{subsec:service-time}

The service rate $\mu$ used in Algorithm~7 is recomputed every
$T_{\mathrm{svc}} = 300\ \mathrm{s}$ by the background measurement
routine described in Equations~\ref{eq:gap}--\ref{eq:svc-blend}. The
pre-deployment calibration in Section~\ref{sec:datagathering}
(Phase~1) supplies the initial seed value for
$\widehat{\mathrm{svc}}_{\mathrm{prev}}$ before any transactions are
logged.

% =====================================================================
\subsection{Ticket Generation}
\label{subsec:ticket}

When a queue number is confirmed, the Ticket Service generates a
printer-agnostic ticket payload comprising four fields: the queue
number, a short alphanumeric code for manual entry, a QR-encoded
status-page URL, and a JSON Web Token (JWT). The JWT is constructed
with the HS256 algorithm:

\begin{equation}
\mathrm{token} \;=\; h \cdot \texttt{.} \cdot p \cdot \texttt{.} \cdot
\mathrm{base64url}\big(\mathrm{HMAC\text{-}SHA256}(h \cdot \texttt{.}
\cdot p,\; K)\big),
\label{eq:jwt}
\end{equation}

\noindent where $h$ is the base64url-encoded header and $p$ is the
base64url-encoded payload containing the queue number (\texttt{sub}),
issue time (\texttt{iat}), expiry (\texttt{exp}), and nonce
(\texttt{jti}); $K$ is the server-side secret. The short code is
drawn from a 32-character unambiguous alphabet (excluding 0, O, 1,
and~I) and formatted as \texttt{XXXX-XXXX}, providing approximately
40 bits of entropy sufficient to resist brute-force guessing within the
four-hour ticket lifetime.

\vspace{0.5em}
\noindent\rule{\linewidth}{0.8pt}

\noindent\textbf{\underline{Algorithm 7: Ticket Generation}}

\noindent\textbf{Input:} \textit{confirmedQueueNumber} $q$,
\textit{statusPageURL} $U$, \textit{signingSecret} $K$\\
\noindent\textbf{Output:} \textit{ticketPayload} (persisted to
database and rendered as PDF or sent to thermal printer)

\begin{enumerate}
  \item \textbf{DRAW} 8 random characters from the unambiguous
        32-character alphabet (excluding 0, O, 1, I);
        \textbf{FORMAT} as \texttt{XXXX-XXXX}
  \item \textbf{GENERATE} $\textit{jti} \gets$ random 32-character
        hexadecimal nonce
  \item \textbf{SET} $\textit{iat} \gets$ current timestamp;
        $\textit{exp} \gets \textit{iat} + 4\ \text{hours}$
  \item \textbf{CONSTRUCT} JWT header $h$ and payload
        $p = \{\mathrm{sub}{:}\,q,\;\mathrm{iat},\;\mathrm{exp},
        \;\mathrm{jti}\}$
  \item \textbf{SIGN} $h.p$ with HMAC-SHA256 and secret $K$ to
        produce \textit{token}
  \item \textbf{ENCODE} QR URL $U(q, \textit{token})$
  \item \textbf{PERSIST} queue number, short code, QR URL,
        \textit{token}, expiry, and status to the MySQL database
  \item \textbf{IF} prototype mode \textbf{THEN RENDER} payload as a
        PDF ticket
  \item \textbf{ELSE SEND} payload to the thermal-printer driver
  \item \textbf{RETURN} ticket payload
\end{enumerate}
\vspace{0.3em}

% =====================================================================
\subsection{Public Display Board with TTS Announcements}
\label{subsec:display-board}

A third presentation surface, alongside the staff dashboard and the
student mobile client, is a public display board mounted in the
service-office waiting area. The board renders the active queue and
announces newly assigned counter calls using the browser's built-in
Web Speech API. All speech synthesis runs client-side; no audio is
transmitted to any third-party service.

A deduplication set indexed by the pair
$(\text{queue\_number}, \text{counter\_number})$ prevents the same
counter assignment from being announced more than once even when the
polling endpoint returns the same active queue across consecutive
intervals. Announcements are spoken sequentially with an $800\ \mathrm{ms}$
gap between utterances to avoid speech overlap when several counter
assignments are issued simultaneously.

% =====================================================================
\subsection{Parameter Summary}
\label{subsec:params}

Table~\ref{tab:all-params} consolidates every tunable parameter
referenced in this section. Values listed are those used in the
deployed prototype.

\begin{table}[!htbp]
\centering
\small
\renewcommand{\arraystretch}{0.94}
\caption{Consolidated parameter values for the deployed \sysname{} prototype.}
\label{tab:all-params}
\begin{tabular}{llcl}
\hline
Subsystem & Symbol & Value & Description \\
\hline
CV Layer        & $A_{\min}$                   & $800\ \mathrm{px}^2$ & Min.\ box area \\
                & $f_{\max}$                   & $0.85$               & Max.\ frame fraction \\
                & $r_{\min}$                   & $0.60$               & Min.\ portrait aspect \\
                & $M_{\min}$                   & $8\ \mathrm{px}$     & Min.\ motion energy \\
                & $c_{\mathrm{byp}}$           & $0.70$               & Static-conf bypass \\
                & $\alpha_{\mathrm{s}}$        & $0.45$               & Bbox EMA factor \\
\hline
Candidate conf. & $N_{\mathrm{conf}}$          & $20$ frames          & Confirmation window \\
                & $\theta_{\mathrm{m}}$        & $1.5\ \mathrm{px}$   & Motion-gate $\sigma$ \\
\hline
Track mgmt.     & $\tau_{\mathrm{iou}}^{\mathrm{rm}}$ & $0.10$   & Remap IoU threshold \\
                & $d_{\max}$                   & $180\ \mathrm{px}$   & Remap distance cap \\
                & $F_{\mathrm{rm}}$            & $45$ frames          & Remap absence window \\
                & $\tau_{\mathrm{iou}}^{\mathrm{dd}}$ & $0.40$   & Dedup IoU threshold \\
                & $f_{\mathrm{dd}}$            & $0.15$               & Centre-fraction gate \\
                & $I_0$                        & $60$ frames          & Re-entry immunity \\
\hline
Face identity   & $\tau_{\mathrm{match}}$      & $0.30$               & Absolute match threshold \\
                & $\tau_{\mathrm{margin}}$     & $0.15$               & Margin over runner-up \\
                & $\tau_{\mathrm{det}}$        & $0.60$               & Face-detector confidence \\
                & $d_{\mathrm{emb}}$           & $512$                & Embedding dimension \\
                & $n_{\mathrm{enrol}}$         & $1$--$5$             & Enrolment photographs \\
                & $N_{\mathrm{face}}$          & $5000$               & System-minted range start \\
\hline
Prediction      & $\alpha_{\mathrm{H}}$        & $0.4$                & Holt's level smoothing \\
                & $\beta_{\mathrm{H}}$         & $0.2$                & Holt's trend smoothing \\
                & $\varphi$                    & $0.85$               & Holt's damping factor \\
                & $h$                          & $20$                 & Holt's projection steps \\
                & $W_{\mathrm{tr}}$            & $20$ samples         & Trend regression window \\
                & $T_{\mathrm{sh}}$            & $600\ \mathrm{s}$    & Short-term horizon \\
                & $W_{\max}^{\text{long}}$     & $60\ \mathrm{min}$   & Long-term cap \\
\hline
Service time    & $T_{\mathrm{svc}}$           & $300\ \mathrm{s}$    & Re-measure interval \\
\hline
Ticket          & N/A                          & $4\ \mathrm{h}$      & JWT expiry \\
                & N/A                          & $40$ bits            & Short-code entropy \\
\hline
Display board   & N/A                          & $800\ \mathrm{ms}$   & TTS inter-utterance gap \\
\hline
\end{tabular}
\end{table}

% =====================================================================
% End of Section 3.2: Design and System Plans
% =====================================================================
